\section{Apimetro}
% ── Frame: Origen de Apimetro ────────────────────────────────────────────────
\begin{frame}{¿Cómo surge Apimetro?}
  \begin{columns}[T]
    \begin{column}{0.58\textwidth}
      \small
      El proyecto Apimetro tiene su génesis entre los años 2023 y 2024,
      derivado de una necesidad cartográfica y de modelación espacial para el universo
      narrativo de \textbf{\enquote{La Saga de la Catacumba}}, proyecto de \termino{World Building}
      que explora la historia de la Ciudad de México.

      \vspace{0.2cm}
      Sin embargo, conforme la historia fue creciendo, se decidió darle un mapa al
      protagonista, el cual contenía las rutas de transporte público y una tarjeta con 4
      viajes gratuitos cada 30 minutos. Para poder modelar el mapa y la tarjeta,
      se necesitaba el panorama completo de la red de transporte, lo cual no existía en un solo lugar.
    \end{column}
    \begin{column}{0.38\textwidth}
      % EDITAR: una línea de explicación por etapa del flujo
      \begin{exampleblock}{\footnotesize 1.\enspace Entrada de datos}
        \scriptsize SHP, GTFS, GeoPackage, CSV, JSON.
      \end{exampleblock}
      % \pause
      \begin{block}{\footnotesize 2.\enspace Procesamiento y limpieza}
        \scriptsize Normalizar todos los datos a un solo formato.
      \end{block}
      % \pause
      \begin{alertblock}{\footnotesize 3.\enspace Salida de datos}
        \scriptsize Se optó por el formato GeoJSON.
      \end{alertblock}
    \end{column}
  \end{columns}
\end{frame}

% ── Frame: Proceso 1 — ETL ───────────────────────────────────────────────────
\begin{frame}{Proceso 1 — ETL: dos fuentes, una base}
  \centering
  \resizebox{0.92\linewidth}{!}{% Diagrama: Proceso 1 de Apimetro — ETL (Extraction, Treatment and Liberation)
% Dos fuentes (datos abiertos + OSM) → ETL en Python (dos fases) → PostgreSQL/PostGIS
% Basado en: TesisUrbanismo 4-1-1-Intro.tex (flujo) y Anexo A2-ETL.tex (fases)
% Usado en: secciones/04-apimetro.tex
\begin{tikzpicture}[
  fuente/.style = {rectangle, rounded corners=5pt,
                   draw=SotmVerde, line width=0.8pt,
                   fill=SotmCrema, text width=4.8cm,
                   align=flush left, minimum height=1.1cm, font=\scriptsize},
  fase/.style   = {rectangle, rounded corners=4pt,
                   draw=SotmVerde!70, fill=white, text width=4.4cm,
                   align=flush left, minimum height=0.8cm, font=\scriptsize},
  marco/.style  = {rectangle, rounded corners=6pt,
                   draw=SotmVerde, line width=0.8pt, dashed,
                   fill=SotmOlivo!12, inner sep=0.25cm},
  result/.style = {rectangle, rounded corners=5pt,
                   fill=SotmTerracota, draw=SotmTerracota, text=white,
                   text width=3.4cm, align=center, minimum height=1.4cm,
                   font=\scriptsize},
  flecha/.style = {-Stealth, thick, color=SotmVerde}
]
  % ── Fuentes ──
  \node[fuente] (f1) at (0,  1.0) {{\small\bfseries Datos abiertos}\\
    GTFS (ADIP) · SHP (INEGI)\\ parámetros operativos (SEMOVI)};
  \node[fuente] (f2) at (0, -1.0) {{\small\bfseries OpenStreetMap}\\
    Mexibús · Mexicable\\ Trolebús Línea 13};

  % ── ETL: dos fases ──
  \node[fase] (p1) at (6.8,  0.5) {\textbf{1.} GTFS → plantilla con campos
    normalizados (líneas, ramales, estaciones)};
  \node[fase] (p2) at (6.8, -0.8) {\textbf{2.} Cruce de geometrías y carga
    a la base de datos};
  \draw[-Stealth, color=SotmVerde!70] (p1) -- (p2);
  \begin{scope}[on background layer]
    \node[marco, fit=(p1)(p2), inner ysep=0.35cm,
          label={[font=\scriptsize\bfseries, text=SotmVerde, yshift=-0.1cm]above:%
          ETL · \textit{Extraction, Treatment and Liberation}}] (etl) {};
  \end{scope}
  \node[font=\tiny\itshape, text=SotmVerde, anchor=north] at (etl.south)
    {Python · GeoPandas};

  % ── Resultado ──
  \node[result] (db) at (12.6, 0) {{\small\bfseries PostgreSQL + PostGIS}\\[2pt]
    estaciones · líneas\\ polígonos};

  % ── Flujo ──
  \draw[flecha] (f1.east) -- (etl.west |- f1.east);
  \draw[flecha] (f2.east) -- (etl.west |- f2.east);
  \draw[flecha] (etl.east) -- (db.west);
\end{tikzpicture}
}\par
  \fuentefigura{Elaboración propia, 2026.}

  \vspace{0.3cm}
  % EDITAR: recordatorio (el detalle de OSM está en el bloque 3)
  {\small\itshape\color{SotmVerde} Como vimos: Mexibús, Mexicable y Trolebús L13 llegan desde OpenStreetMap.}
\end{frame}

% ── Frame: E-R ───────────────────────────────────────────────────
\begin{frame}{Esquema E-R de Apimetro}
  \begin{columns}[c]
    \begin{column}{0.45\textwidth}
      \centering
      \resizebox{\linewidth}{!}{% Diagrama: Esquema Entidad-Relación simplificado de Apimetro
% Tres entidades: Estaciones (nodos puntuales), Líneas (trazos con métricas),
%   Polígonos (áreas administrativas). Sin columnas individuales.
% Usado en: 4-Capas-Codigo/4-1-Apimetro/4-1-2-Modelo-Datos.tex  (fig:apimetro-er-simple)
\tikzsetnextfilename{apimetro-er-simple}
\begin{tikzpicture}[
  entidad/.style={
    rectangle, draw=unamAzul, fill=unamAzul!12,
    thick, rounded corners=4pt,
    minimum width=3.4cm, minimum height=1.1cm,
    align=center, font=\small\bfseries
  },
  relacion/.style={Stealth-Stealth, thick, color=unamAzul!60},
  card/.style={font=\footnotesize\color{unamAzul!80}},
  etiqrel/.style={
    font=\footnotesize\itshape\color{gray!65},
    fill=white, inner sep=2pt
  }
]

%% — Vértices del triángulo —
\node[entidad] (est) at (0,    0)    {Estaciones};
\node[entidad] (lin) at (6.5,  0)    {Líneas};
\node[entidad] (pol) at (3.25, -3.1) {Polígonos};

%% — Cardinalidades ancladas a cada flecha con pos= (no coordenadas fijas):
%%    así no se enciman con las cajas si cambia la fuente o el tamaño —
%% — Estaciones ↔ Líneas (N:M) —
\draw[relacion] (est) -- node[etiqrel, above] {agrupa / conecta}
  node[card, pos=0.08, below] {N}
  node[card, pos=0.92, below] {M} (lin);

%% — Estaciones ↔ Polígonos (N:1) —
\draw[relacion] (est) -- node[etiqrel, left, xshift=-4pt] {ubicada en}
  node[card, pos=0.12, right, xshift=2pt] {N}
  node[card, pos=0.88, right, xshift=2pt] {1} (pol);

%% — Líneas ↔ Polígonos (M:N) —
\draw[relacion] (lin) -- node[etiqrel, right, xshift=4pt] {atraviesa}
  node[card, pos=0.12, left, xshift=-2pt] {M}
  node[card, pos=0.88, left, xshift=-2pt] {N} (pol);

\end{tikzpicture}
}
      \fuentefigura{Elaboración propia, 2026.}
    \end{column}
    \begin{column}{0.53\textwidth}
      \footnotesize
      \vspace*{-0.3cm}
      \begin{itemize}\setlength{\itemsep}{0.05cm}
        \item \textbf{Estaciones:} nodos puntuales de la red con atributos de
              conectividad, sistema de transporte y ubicación administrativa.
        \item \textbf{Líneas:} trazos de los ramales con las métricas operativas
              necesarias para el cálculo posterior (velocidad promedio, frecuencia de paso,
              capacidad del vehículo).
        \item \textbf{Polígonos:} áreas administrativas para filtrar cualquier
              consulta por alcaldía o municipio.
        \item \textbf{Relaciones:} las líneas se vinculan a las estaciones mediante
              un identificador de sistema y ramal, lo que permite reconstruir el grafo
              metropolitano.
      \end{itemize}
    \end{column}
  \end{columns}
\end{frame}

% ── Frame: Proceso 2 — Salida GeoJSON ────────────────────────────────────────
% Comillas JSON con \textquotedbl: el " está activo con babel-spanish
\providecommand{\jq}[1]{\textquotedbl#1\textquotedbl}
\begin{frame}[t]{Proceso 2 — Salida: API REST en GeoJSON}
  \begin{columns}[T]
    \begin{column}{0.47\textwidth}
      \centering
      \resizebox{\linewidth}{!}{% Diagrama: Dos tipos de salida de Apimetro
% Columna izquierda: JSON Descriptivo (atributos alfanuméricos)
% Columna derecha: GeoJSON Geográfico (geometría + métricas)
% Usado en: 4-Capas-Codigo/4-1-Apimetro/4-1-4-Salidas.tex (fig:apimetro-salidas-tipos)
% Copia para la presentación: coordinates y tipo de geometría en dos renglones
\tikzsetnextfilename{apimetro-salidas-tipos}
\begin{tikzpicture}[
  cabecera/.style={
    rectangle, draw=unamAzul, fill=unamAzul, text=white,
    minimum width=5.4cm, minimum height=0.9cm,
    align=center, font=\small\bfseries
  },
  cuerpo/.style={
    rectangle, draw=unamAzul!40, fill=unamAzul!8,
    minimum width=5.4cm, text width=5.0cm,
    align=left, inner sep=8pt, font=\footnotesize
  },
  uso/.style={
    rectangle, draw=unamOro!70, fill=unamOro!18,
    minimum width=5.4cm, minimum height=0.75cm,
    align=center, font=\footnotesize\itshape, text=unamAzul
  }
]

%% ── Columna izquierda: JSON Descriptivo ──────────────────────────
\node[cabecera] (jh) at (0, 0) {JSON Descriptivo};

\node[cuerpo, anchor=north] (jb) at (jh.south) {%
  \textbf{¿Qué contiene?}\\[4pt]
  \texttt{linea\_id} · \texttt{nombre}\\
  \texttt{sistema} · \texttt{color\_esp}\\
  \texttt{clasificacion}\\[4pt]
  \textit{Sin geometría}
};

\node[uso, anchor=north] (ju) at (jb.south)
  {tablas · fichas · simbología};

%% ── Columna derecha: GeoJSON Geográfico ──────────────────────────
\node[cabecera] (gh) at (6.6, 0) {GeoJSON Geográfico};

\node[cuerpo, anchor=north] (gb) at (gh.south) {%
  \textbf{¿Qué contiene?}\\[4pt]
  \texttt{coordinates}\\
  \quad Point / MultiLineString\\
  \texttt{properties}\\
  \quad velocidad · frecuencia · capacidad\\[4pt]
  \textit{Listo para web, sin conversión}
};

\node[uso, anchor=north] (gu) at (gb.south)
  {mapas · ruteo · \textit{VFTModel}};

\end{tikzpicture}
}
      \fuentefigura{Elaboración propia, 2026.}

      \raggedright\scriptsize
      \vspace*{-0.2cm}
      % EDITAR: porqués del GeoJSON y del modelo SaaS
      \begin{itemize}\setlength{\itemsep}{0.05cm}
        \item \textbf{Se procesa una vez, se consulta siempre:} cualquier cliente
              recibe la red ya normalizada desde \texttt{apimetro.dev}.
        \item \textbf{Texto plano:} legible, versionable y nativo en la web.
        \item \textbf{Una línea de Python} lo convierte a GeoPackage para QGIS o cualquier otro software cartográfico.
      \end{itemize}
    \end{column}
    \begin{column}{0.51\textwidth}
      % EDITAR: fragmento reducido de una estación (Anexo A6 de la tesis)
      \begin{block}{\footnotesize ¿Qué devuelve una consulta? Una estación}
        \scriptsize
        \begin{tabular}{@{}l@{\enspace}l@{}}
          \ttfamily \{ \jq{geometry}: \{ & \\
          \ttfamily\quad \jq{type}: \jq{Point}, &
            {\color{SotmTerracota}← \itshape es un punto} \\
          \ttfamily\quad \jq{coordinates}: & \\
          \ttfamily\quad\quad [-99.134, 19.429] \}, &
            {\color{SotmTerracota}← \itshape dónde está} \\
          \ttfamily\ \ \jq{properties}: \{ & \\
          \ttfamily\quad \jq{nombre}: \jq{20 de Noviembre}, &
            {\color{SotmTerracota}← \itshape qué es} \\
          \ttfamily\quad \jq{sistema}: \jq{MB}, &
            {\color{SotmTerracota}← \itshape a qué red} \\
          \ttfamily\quad \jq{alcaldia}: \jq{Cuauhtémoc}, &
            {\color{SotmTerracota}← \itshape para filtrar} \\
          \ttfamily\quad \jq{es\_cetram}: false \} \} &
            {\color{SotmTerracota}← \itshape ¿es CETRAM?} \\
        \end{tabular}
      \end{block}
      \vspace{0.1cm}
      \footnotesize Las líneas viajan igual, como \texttt{MultiLineString}, con
      velocidad, frecuencia y capacidad en sus \texttt{properties}: listas para
      calcular sin cruzar archivos.
    \end{column}
  \end{columns}
\end{frame}

% ── Frame: Sistemas integrados ───────────────────────────────────────────────
% Cifras de Apimetro/db/init/04_seed.sql (26 oct 2026): ubicaciones únicas por sistema
%   (22 878 registros → 11 115 ubicaciones; la suma por sistema da 11 154 porque
%   39 paradas son compartidas entre sistemas). Coincide con los nodos del grafo (tesis 5.2)
\begin{frame}{¿Qué transportes integra Apimetro?}
  \begin{columns}[c]
    \begin{column}{0.58\textwidth}
      \centering\scriptsize
      \begin{tabular}{@{}lrr@{}}
        \toprule
        \textbf{Sistema} & \textbf{Líneas o rutas} & \textbf{Estaciones y paradas} \\
        \midrule
        Metro                    &   12 &     195 \\
        Metrobús                 &    8 &     373 \\
        Tren Ligero              &    1 &      18 \\
        Tren Suburbano           &    1 &       7 \\
        Tren Interurbano         &    1 &       5 \\
        Cablebús                 &    3 &      19 \\
        Mexicable                &    2 &      14 \\
        Mexibús                  &    6 &     157 \\
        Trolebús                 &   12 &     740 \\
        RTP                      &  114 & 4\,907 \\
        Corredores Concesionados &  137 & 4\,629 \\
        Pumabús                  &   12 &      90 \\
        \midrule
        \textbf{Total}           & \textbf{309} & \textbf{11\,115}\textsuperscript{*} \\
        \bottomrule
      \end{tabular}

      {\tiny\textsuperscript{*}\,Cada parada se cuenta una sola vez, aunque la usen varias rutas.}

      \fuentefigura{Elaboración propia con datos de Apimetro, 2026.}
    \end{column}
    \begin{column}{0.38\textwidth}
      \centering
      {\Huge\bfseries\color{SotmVerde} 12}\\[-0.05cm]
      {\small sistemas de transporte}\\[0.25cm]
      {\Huge\bfseries\color{SotmVerde} 309}\\[-0.05cm]
      {\small líneas y rutas}\\[0.25cm]
      {\Huge\bfseries\color{SotmTerracota} 11\,115}\\[-0.05cm]
      {\small estaciones y paradas}\\[0.3cm]
      % EDITAR: frase de cierre
      {\footnotesize\itshape De la CDMX y del Estado de México, en una sola base.}
    \end{column}
  \end{columns}
\end{frame}

% ── Frame: Líneas por sistema ────────────────────────────────────────────────
\begin{frame}{Las líneas, una por una}
  \centering\scriptsize
  \begin{tabular}{@{}lp{0.68\textwidth}@{}}
    \toprule
    \textbf{Sistema} & \textbf{Líneas} \\
    \midrule
    Metro            & 1, 2, 3, 4, 5, 6, 7, 8, 9, A, B y 12 \\
    Metrobús         & 1, 2, 3, 4, 5, 6, 7 y SL01 \\
    Tren Ligero      & 1 línea \\
    Tren Suburbano   & 1 línea \\
    Tren Interurbano & El Insurgente \\
    Cablebús         & 1, 2 y 3 \\
    Mexicable        & 1 y 2 \\
    Mexibús          & 01, 1, 2, 2A, 3 y 4 \\
    Trolebús         & 1 a 10, 12 y 13 \\
    \midrule
    RTP              & 114 rutas \enspace{\itshape(p.\,ej. Ruta 1-D: Metro Santa Martha – Metro Mixcoac)} \\
    Corredores Conc. & 137 rutas \enspace{\itshape(p.\,ej. Ruta 10C: RCA Víctor – Peñón)} \\
    Pumabús          & 12 rutas \enspace{\itshape(p.\,ej. PUMA1: Metro Universidad – Islas)} \\
    \bottomrule
  \end{tabular}

  \fuentefigura{Elaboración propia con datos de Apimetro, 2026.}
\end{frame}

% ── Frame: Homologación ──────────────────────────────────────────────────────
\begin{frame}{Un solo estándar, sin perder detalles}
  \begin{columns}[T]
    \begin{column}{0.48\textwidth}
      % EDITAR: el problema de la homologación
      \begin{alertblock}{\footnotesize El problema: cada sistema, su propio formato}
        \scriptsize
        \begin{itemize}\setlength{\itemsep}{0.05cm}
          \item Cada transporte tiene su organismo, su sindicato o concesionario
                y sus particularidades.
          \item Cada uno registra sus datos a su manera: nombres, claves,
                sentidos de ida y regreso, formatos distintos.
          \item El dato existe, pero no se puede cruzar: falta
                \textbf{homologación}.
        \end{itemize}
      \end{alertblock}
    \end{column}
    \begin{column}{0.48\textwidth}
      % EDITAR: lo que resuelve la normalización
      \begin{block}{\footnotesize La solución: un mismo formato para todos}
        \scriptsize
        \begin{itemize}\setlength{\itemsep}{0.05cm}
          \item Toda línea y estación se describe con los mismos campos:
                nombre, sistema, color oficial, año, alcaldía o municipio.
          \item Sin omitir lo particular: velocidad, frecuencia, capacidad,
                tipo de vía y si la estación es CETRAM.
          \item Resultado: una base estándar que se puede usar de forma
                \textbf{independiente}, sin depender de cada operador.
        \end{itemize}
      \end{block}
    \end{column}
  \end{columns}

  \vspace{0.4cm}
  \centering
  {\bfseries\color{SotmVerde} 12 sistemas · 12 formatos} \enspace
  {\color{SotmTerracota}\Large$\rightarrow$} \enspace
  {\bfseries\color{SotmTerracota} 1 estándar}
\end{frame}

% ── Frame: Acceso a los datos ────────────────────────────────────────────────
\begin{frame}{Tres formas de usar los datos}
  \begin{columns}[T]
    \begin{column}{0.31\textwidth}
      % EDITAR
      \begin{exampleblock}{\footnotesize 1.\enspace En la web}
        \scriptsize
        \texttt{apimetro.dev}

        \vspace{0.1cm}
        Consulta desde el navegador con la documentación interactiva
        (Swagger), sin instalar nada.

        \vspace{0.1cm}
        {\itshape Para cualquier persona.}
      \end{exampleblock}
    \end{column}
    \begin{column}{0.31\textwidth}
      % EDITAR
      \begin{block}{\footnotesize 2.\enspace La base de datos}
        \scriptsize
        \texttt{seed.sql} ($\approx$\,38 MB)

        \vspace{0.1cm}
        Se genera con el ETL y carga la red completa en PostgreSQL + PostGIS,
        lista para QGIS. \textbf{Próximamente en GitHub.}

        \vspace{0.1cm}
        {\itshape Para analistas SIG.}
      \end{block}
    \end{column}
    \begin{column}{0.31\textwidth}
      % EDITAR
      \begin{alertblock}{\footnotesize 3.\enspace El código completo}
        \scriptsize
        \texttt{github.com/}\\
        \texttt{galigaribaldi/Apimetro}

        \vspace{0.1cm}
        Se descarga y se levanta en cualquier equipo con Docker Compose.

        \vspace{0.1cm}
        {\itshape Para desarrolladores.}
      \end{alertblock}
    \end{column}
  \end{columns}

  \vspace{0.4cm}
  \centering\scriptsize
  Licencia \textit{Business Source License} 1.1: uso gratuito académico,\\
  de investigación y del sector público · Apache 2.0 a partir de 2029
\end{frame}
